% paper_mon(7).pdf болон study_three_portfolios_rolling_HAC_with_EW_costs(3).py-д
% тулгуурласан нэгтгэсэн арга зүйн тайлбар.
% ЭНЭ НЬ 2.4, 2.5-Р ХЭСГИЙН ТОВЧИЛСОН ОРЛУУЛАХ ХУВИЛБАР.
% Өмнөх 2.4.1-ийн Лемм 1, Теорем 1 болон баталгааг хамтад нь хасна.

\section{Ковариацын матрицын агшилтын үнэлгээ}

\subsection{Нийтлэг тэмдэглэгээ ба шугаман агшилт}
$n$ сарын өгөөж, $N$ хөрөнгийн хувьд $y_{ti}$-г $t$-р сарын
$i$-р хөрөнгийн өгөөж, $\bar y_i=n^{-1}\sum_{t=1}^n y_{ti}$-г
түүврийн дундаж гэе. $x_{ti}=y_{ti}-\bar y_i$ болон
$X=(x_{ti})$ гэж тэмдэглэвэл энэхүү судалгааны Python
хэрэгжүүлэлтэд
\[
 S=\frac{1}{n}X^\top X,\qquad
 s_{ij}=\frac{1}{n}\sum_{t=1}^n x_{ti}x_{tj}
\]
гэж тооцно. Энэ нь $1/(n-1)$ хуваарьтай хазайлтгүй
түүврийн ковариацын үнэлэгчээс ялгаатай. Бие даасан, ижил
тархалттай ажиглалтын хувьд $E[S]=(n-1)\Sigma/n$ байна.

Шугаман агшилтын үнэлэгчийн ерөнхий хэлбэр
\[
 \widehat\Sigma_k=(1-\widehat\alpha_k)S+
 \widehat\alpha_k T_k,\qquad 0\le\widehat\alpha_k\le1,
\]
энд $T_k$ нь тухайн аргын зорилтот матриц юм.
Тооцооллын нийтлэг хэмжигдэхүүнүүдийг
\[
 z_{tij}=x_{ti}x_{tj},\qquad
 \widehat b=\frac{1}{n^2N}\sum_{t=1}^n\sum_{i=1}^N
 \sum_{j=1}^N(z_{tij}-s_{ij})^2,
 \qquad
 \widehat d_k=\frac1N\sum_{i=1}^N\sum_{j=1}^N
 (s_{ij}-(T_k)_{ij})^2
\]
гэж авна. $\widehat d_k>0$ үед доор өгсөн харьцааг
$[0,1]$ завсарт хязгаарлаж агшилтын эрчмийг тооцно.
Эдгээр нь кодод хэрэгжүүлсэн өгөгдөлд суурилсан
коэффициентийн үнэлгээнүүд юм.

Зорилтот матрицууд $S$-ээс тооцогддог тул санамсаргүй
бус $T$ болон $E[S]=\Sigma$ нөхцөлд гаргасан
$\alpha^*=b/d$ гэсэн төгсгөлөг түүврийн баталгааг эдгээр
үнэлэгчид шууд хэрэглэхгүй. Иймээс энд асимптотик
аргачлалын хэрэгжүүлэлтийг тодорхойлж байгаа бөгөөд
төслийн өгөгдөлд яг оновчтой болохыг батлаагүй.

\subsection{Тогтмол вариацын зорилтот матриц (Var)}
\[
 \widehat\mu=\frac{\operatorname{tr}(S)}{N},\qquad
 T_{\mathrm{var}}=\widehat\mu I_N,\qquad
 \widehat\alpha_{\mathrm{var}}=
 \left[\frac{\widehat b}{\widehat d_{\mathrm{var}}}\right]_0^1.
\]
Энд $[a]_0^1=\min\{1,\max\{0,a\}\}$.
$\widehat d_{\mathrm{var}}=0$ үед код
$\widehat\alpha_{\mathrm{var}}=1$ гэж авна.

\subsection{Диагональ зорилтот матриц (Diag)}
\[
 T_{\mathrm{diag}}=\operatorname{diag}(s_{11},\ldots,s_{NN}),
 \qquad
 \widehat b_{\mathrm{diag}}=\frac{1}{n^2N}
 \sum_{t=1}^n\sum_{i\ne j}(z_{tij}-s_{ij})^2,
\]
\[
 \widehat\alpha_{\mathrm{diag}}=
 \left[\frac{\widehat b_{\mathrm{diag}}}
 {\widehat d_{\mathrm{diag}}}\right]_0^1.
\]
Энэ арга зөвхөн диагоналийн бус элементүүдийг агшаана.
$\widehat d_{\mathrm{diag}}=0$ үед код коэффициентийг
нэгтэй тэнцүүлнэ.

\subsection{Тогтмол корреляцын зорилтот матриц (Corr)}
\[
 r_{ij}=\frac{s_{ij}}{\sqrt{s_{ii}s_{jj}}},\qquad
 \bar r=\frac{1}{N(N-1)}\sum_{i\ne j}r_{ij},
\]
\[
 (T_{\mathrm{corr}})_{ij}=
 \begin{cases}s_{ii},&i=j,\\
 \bar r\sqrt{s_{ii}s_{jj}},&i\ne j.
 \end{cases}
\]
Кодод $v_i=n^{-1}\sum_t(x_{ti}^2-s_{ii})^2$ болон
\[
 \theta_{ii,ij}=\frac1n\sum_t(x_{ti}^2-s_{ii})
 (x_{ti}x_{tj}-s_{ij}),\qquad
 \theta_{jj,ij}=\frac1n\sum_t(x_{tj}^2-s_{jj})
 (x_{ti}x_{tj}-s_{ij})
\]
гэж тооцоод
\[
 \widehat c_{\mathrm{corr}}=\frac{1}{nN}
 \left[\sum_{i=1}^N v_i+
 \frac{\bar r}{2}\sum_{i\ne j}
 \left(\sqrt{\frac{s_{jj}}{s_{ii}}}\theta_{ii,ij}+
 \sqrt{\frac{s_{ii}}{s_{jj}}}\theta_{jj,ij}\right)\right],
\]
\[
 \widehat\alpha_{\mathrm{corr}}=
 \left[\frac{\widehat b-\widehat c_{\mathrm{corr}}}
 {\widehat d_{\mathrm{corr}}}\right]_0^1
\]
гэж ашиглана. $\widehat d_{\mathrm{corr}}=0$ үед
коэффициентийг нэгтэй тэнцүүлнэ.

\subsection{Нэг хүчин зүйлийн зорилтот матриц (SF)}
Зах зээлийн хүчин зүйлийг S\&P 500 индексийн сарын өгөөж
$m_t$-ээр төлөөлүүлнэ. $u_t=m_t-\bar m$ гэвэл
\[
 \widehat\sigma_m^2=\frac1n\sum_tu_t^2,\quad
 \widehat\beta_i=\frac{\sum_t x_{ti}u_t}{\sum_tu_t^2},\quad
 e_{ti}=x_{ti}-\widehat\beta_i u_t,\quad
 \widehat\sigma_{\varepsilon,i}^2=\frac1n\sum_te_{ti}^2.
\]
\[
 T_{\mathrm{SF}}=\widehat\sigma_m^2
 \widehat\beta\widehat\beta^\top+
 \operatorname{diag}(\widehat\sigma_{\varepsilon,1}^2,
 \ldots,\widehat\sigma_{\varepsilon,N}^2).
\]
Кодын хам хэлбэлзлийн хэмжигдэхүүнийг
\[
 a_{tij}=\begin{cases}
 x_{ti}^2-s_{ii},&i=j,\\
 \widehat\beta_j x_{ti}u_t+
 \widehat\beta_i x_{tj}u_t-
 \widehat\beta_i\widehat\beta_j u_t^2-
 (T_{\mathrm{SF}})_{ij},&i\ne j
 \end{cases}
\]
гэж тодорхойлж,
\[
 \widehat c_{\mathrm{SF}}=\frac{1}{nN}
 \sum_{i=1}^N\sum_{j=1}^N
 \left\{\frac1n\sum_{t=1}^n
 a_{tij}(z_{tij}-s_{ij})\right\},
 \qquad
 \widehat\alpha_{\mathrm{SF}}=
 \left[\frac{\widehat b-\widehat c_{\mathrm{SF}}}
 {\widehat d_{\mathrm{SF}}}\right]_0^1
\]
гэж авна. $\widehat d_{\mathrm{SF}}=0$ үед
коэффициентийг нэгтэй тэнцүүлнэ.

\subsection{Аналитик шугаман бус агшилт (NLS)}
$S=U\Lambda U^\top$, $\Lambda=\operatorname{diag}
(\lambda_1,\ldots,\lambda_N)$ гэе. Кодод
$c=N/n$, $h=n^{-1/3}$, $h_j=\lambda_jh$ гэж авч,
Epanechnikov kernel $k$ болон түүний Hilbert
хувиргалт $Hk$-д тулгуурлан
\[
 \widehat f(\lambda_i)=\frac1N\sum_{j=1}^N
 \frac{1}{h_j}k\!\left(\frac{\lambda_i-\lambda_j}{h_j}\right),
\]
\[
 \widehat{Hf}(\lambda_i)=\frac1N\sum_{j=1}^N
 \frac{1}{h_j}Hk\!\left(\frac{\lambda_i-\lambda_j}{h_j}\right)
\]
гэж тооцно. Үүний дараа
\[
 \widehat d_i=\frac{\lambda_i}{
 [1-c-c\pi\lambda_i\widehat{Hf}(\lambda_i)]^2+
 [c\pi\lambda_i\widehat f(\lambda_i)]^2},\qquad
 \widehat\Sigma_{\mathrm{NLS}}=
 U\operatorname{diag}(\widehat d_1,\ldots,\widehat d_N)U^\top.
\]
Кодод тоон тогтвортой байдлын үүднээс $\lambda_i$-г
доод тал нь $10^{-8}$ байхаар хязгаарладаг.
Энэхүү $N=10$, $n=120$ нөхцөлд асимптотик
NLS-ийн ойролцооллын чанарыг зөвхөн томьёогоор
баталгаажуулах боломжгүй; тоон баталгаажуулалт
шаардлагатай.

% ЧУХАЛ: Эндэх томьёонууд нь өгсөн Python кодын
% хэрэгжүүлэлтийг тайлбарлана. Тэдгээрийн бүх асимптотик
% онолын нөхцөлийг бие даан баталсан гэсэн үг биш.
